\begin{pastexam}[label={px:ntu-112-dsa-15}]{臺大資工所 112 學年度 第 15 題}
（問答；只看最後答案，完全正確才給分）

\textbf{(21 points)} Given $k$ singly linked lists, each of which has $n$ nodes. The numbers in the
nodes of the $i$-th list are given by $a_{i,1},a_{i,2},\ldots,a_{i,n}$, as shown in the figure below.
Each of the $k$ lists has the numbers sorted in non-decreasing order, \textit{i.e.},
$a_{i,1}\le a_{i,2}\le\cdots\le a_{i,n}$, where $i$ is the index of the list. Professor Q asks the
students to develop an algorithm to merge these k lists into one singly linked list sorted in
non-decreasing order. Three students have come up with different algorithms, which are given below
in pseudo code.

\begin{center}
\begin{tikzpicture}[font=\sffamily\small,>={Stealth[length=1.8mm]},c/.style={draw,minimum width=6.5mm,minimum height=6mm,inner sep=0pt,font=\scriptsize}]
\node[font=\sffamily\bfseries,anchor=east] at (-0.6,0.00) {list 1};
\node[c] (d00) at (0.00,0.00) {$a_{1,1}$};
\node[c,right=-\pgflinewidth of d00] (p00) {};
\node[c] (d01) at (1.70,0.00) {$a_{1,2}$};
\node[c,right=-\pgflinewidth of d01] (p01) {};
\node[c] (d02) at (3.40,0.00) {$a_{1,3}$};
\node[c,right=-\pgflinewidth of d02] (p02) {};
\node[c] (d03) at (6.00,0.00) {$a_{1,n}$};
\node[c,right=-\pgflinewidth of d03] (p03) {};
\node[font=\sffamily\bfseries] (h0) at (0,0.95) {head};
\draw[->] (h0.south)--(d00.north);
\draw[->] (p00.center)--(d01.west);
\draw[->] (p01.center)--(d02.west);
\draw[->] (p02.center)--++(0.55,0);
\node at (4.98,0.00) {$\cdots$};
\draw[->] (5.32,0.00)--(d03.west);
\draw (p03.south west)--(p03.north east);
\node[font=\sffamily\bfseries,anchor=east] at (-0.6,-1.90) {list 2};
\node[c] (d10) at (0.00,-1.90) {$a_{2,1}$};
\node[c,right=-\pgflinewidth of d10] (p10) {};
\node[c] (d11) at (1.70,-1.90) {$a_{2,2}$};
\node[c,right=-\pgflinewidth of d11] (p11) {};
\node[c] (d12) at (3.40,-1.90) {$a_{2,3}$};
\node[c,right=-\pgflinewidth of d12] (p12) {};
\node[c] (d13) at (6.00,-1.90) {$a_{2,n}$};
\node[c,right=-\pgflinewidth of d13] (p13) {};
\node[font=\sffamily\bfseries] (h1) at (0,-0.95) {head};
\draw[->] (h1.south)--(d10.north);
\draw[->] (p10.center)--(d11.west);
\draw[->] (p11.center)--(d12.west);
\draw[->] (p12.center)--++(0.55,0);
\node at (4.98,-1.90) {$\cdots$};
\draw[->] (5.32,-1.90)--(d13.west);
\draw (p13.south west)--(p13.north east);
\node[font=\sffamily\bfseries,anchor=east] at (-0.6,-4.30) {list k};
\node[c] (d20) at (0.00,-4.30) {$a_{k,1}$};
\node[c,right=-\pgflinewidth of d20] (p20) {};
\node[c] (d21) at (1.70,-4.30) {$a_{k,2}$};
\node[c,right=-\pgflinewidth of d21] (p21) {};
\node[c] (d22) at (3.40,-4.30) {$a_{k,3}$};
\node[c,right=-\pgflinewidth of d22] (p22) {};
\node[c] (d23) at (6.00,-4.30) {$a_{k,n}$};
\node[c,right=-\pgflinewidth of d23] (p23) {};
\node[font=\sffamily\bfseries] (h2) at (0,-3.35) {head};
\draw[->] (h2.south)--(d20.north);
\draw[->] (p20.center)--(d21.west);
\draw[->] (p21.center)--(d22.west);
\draw[->] (p22.center)--++(0.55,0);
\node at (4.98,-4.30) {$\cdots$};
\draw[->] (5.32,-4.30)--(d23.west);
\draw (p23.south west)--(p23.north east);
\node[font=\sffamily\bfseries] at (-1.0,-3.0) {$\cdots$};
\end{tikzpicture}
\end{center}

\noindent Student A:\\[2pt]
\begin{tabular}{@{\hspace{1.5em}}l@{\hspace{1em}}p{0.8\linewidth}@{}}
L1. & Create an empty singly linked list $S$ to collect the result. \\
L2. & DO \\
L3. & \quad Iterate over the numbers in the head nodes of the $k$ lists, and find the node with
the smallest number, denoted as node $M$. \\
L4. & \quad Remove node $M$ from the original list, and insert it into the result list $S$ from the
tail. \\
L5. & UNTIL all original $k$ lists are empty \\
L6. & Return the result list $S$, which contains all nodes from the original $k$ lists and sorted
in non-decreasing order.
\end{tabular}

\medskip
\noindent Student B:\\[2pt]
\begin{tabular}{@{\hspace{1.5em}}l@{\hspace{1em}}p{0.8\linewidth}@{}}
L1. & Create an empty singly linked list $S$ to collect the result. \\
L2. & Create an empty min heap $H$. \\
L3. & Remove the $k$ head nodes from the $k$ lists and insert them into $H$. Each node in $H$ also
contains the index number $i$ of the list where it is from. \\
L4. & DO \\
L5. & \quad Remove the min node $M$ from the heap. Insert this node into $S$ from the tail. Take note
of the index number stored in $M$, denoted as $m$. \\
L6. & \quad If list $m$ is not empty, remove the head node from list $m$ and insert it into $H$. \\
L7. & WHILE $H$ is not empty \\
L8. & Return the result list $S$, which contains all nodes from the original $k$ lists and sorted
in non-decreasing order.
\end{tabular}

\medskip
\noindent Student C:\\
Call the recursive divide-and-conquer function MergeTwoGroupLists (\textit{k sorted linked lists}).
The return value of the function is a list containing all nodes from the original $k$ lists and
sorted in non-decreasing order.

\smallskip
\begin{tabular}{@{\hspace{1.5em}}l@{\hspace{1em}}p{0.8\linewidth}@{}}
L1. & Function MergeTwoGroupLists ($m$ sorted linked lists) \\
L2. & \quad If $m{=}{=}1$ then return the only list from the input. \\
L3. & \quad Divide the input $m$ lists into two groups of lists, with $\lceil \frac{m}{2}\rceil$ and
$\lfloor \frac{m}{2}\rfloor$ lists, respectively. \\
L4. & \quad Recursively call MergeTwoGroupLists for each of these two groups of lists, merging the
first group of the lists into one sorted list, $S_1$, and the second group of the lists into one
sorted list, $S_2$. \\
L5. & \quad Merge the two sorted lists $S_1$ and $S_2$ into one sorted list $S$. \\
L6. & \quad Return $S$.
\end{tabular}

\medskip
For each of these three algorithms, analyze and give their asymptotic worst-case running times with
\textbf{the big-O notation} and \textbf{in terms of $k$ and $n$}.

Note:
\begin{itemize}
\item Lower-order terms and constant factors (excluding $k$ and $n$) should be removed in the big-O
notation in the answer.
\item The bound in the answer needs to be tight.
\item Only the final result will be graded and only fully correct answers will be given points.
Please clearly mark your final result in the answer sheet.
\end{itemize}
\end{pastexam}

\begin{pastsoln}
答案：學生 A $O(nk^2)$；學生 B $O(nk \log k)$；學生 C $O(nk \log k)$。

三個演算法都是 $k$ 路合併，總共要輸出 $N = nk$ 個節點。 搬移一個節點（從原串列拿下、接到 $S$ 的
尾端）只改幾個指標，$O(1)$；所以差別在「每次怎麼找出最小的那一個」。 以下假設 $k \ge 2$（$k = 1$
時三個都是 $O(n)$）。

\textbf{學生 A：$O(nk^2)$。} L3 每一輪要看 $k$ 條串列的開頭，$\Theta(k)$；L4 搬一個節點 $O(1)$。
每一輪恰好輸出一個節點，所以共 $nk$ 輪，總時間 $nk \cdot \Theta(k) = \Theta(nk^2)$。 這是緊的：
即使把已經空了的串列跳過，最壞情況也可以讓所有串列一起變短（例如第 $i$ 條串列是
$i, i + k, i + 2k, \ldots$），前 $nk - k$ 輪每條串列都還有節點，每輪都真的要比 $k$ 個開頭。

\textbf{學生 B：$O(nk \log k)$。} 這是資結書第~8.3~節「合併 $k$ 個排好的數列」的做法。
\begin{itemize}
\item L3：$k$ 次插入，$O(k \lg k)$（或用建堆積 $O(k)$）。
\item L4--L7 的迴圈：每一輪取出一個節點輸出，所以共 $nk$ 輪。 每輪一次 \proc{Extract-Min}，最多一次
插入。 堆積裡每條串列最多同時有一個節點，大小 $\le k$，每次操作 $O(\lg k)$。
\end{itemize}
總時間 $O(k \lg k) + nk \cdot O(\lg k) = O(nk \lg k)$。 緊的理由和 A 一樣：同樣的輸入下大部分時間
堆積裡都有 $k$ 個節點，二元堆積每次操作 $\Theta(\lg k)$。

\textbf{學生 C：$O(nk \log k)$。} 合併兩條總長 $L$ 的排好的串列（L5）是 $\Theta(L)$。 設 $T(m)$ 是
合併 $m$ 條長 $n$ 的串列的時間：
\[
T(m) = T(\lceil m/2 \rceil) + T(\lfloor m/2 \rfloor) + \Theta(mn), \qquad T(1) = \Theta(1).
\]
用遞迴樹（資結書第~2~章）：第 $j$ 層把 $k$ 條串列分成 $2^j$ 組，每一組合併的總長加起來還是全部的
$nk$ 個節點，所以每一層的合併共 $\Theta(nk)$。 每一層組數加倍，到每組只剩一條串列為止，共
$\lceil \lg k \rceil$ 層。 總時間 $\Theta(nk) \cdot \lceil \lg k \rceil = \Theta(nk \lg k)$。 （另外每個
遞迴呼叫本身 $O(1)$，共 $2k - 1$ 次呼叫，$O(k)$，被吸收。）

三者的 big-O 寫法：A 是 $O(nk^2)$，B、C 都是 $O(nk \log k)$。 B 和 C 一樣快，而且都比 A 好
$k / \lg k$ 倍。
\end{pastsoln}
